\documentclass[12pt,a4paper]{article}
\usepackage[UTF8]{ctex}
\usepackage{geometry}
\usepackage{graphicx}
\usepackage{listings}
\usepackage{xcolor}
\usepackage{booktabs}
\usepackage{amsmath}
\usepackage{amssymb}
\usepackage{float}
\usepackage{hyperref}
\usepackage{titlesec}
\usepackage{setspace}
\usepackage{indentfirst}

\geometry{left=2.5cm,right=2.5cm,top=2.5cm,bottom=2.5cm}

\setlength{\parindent}{2em}

\titleformat{\section}{\centering\heiti\fontsize{16pt}{\baselineskip}\selectfont\bfseries}{第\chinese{section}章}{1em}{}
\titlespacing{\section}{0pt}{24pt}{18pt}

\titleformat{\subsection}{\centering\heiti\fontsize{14pt}{\baselineskip}\selectfont\bfseries}{第\chinese{subsection}节}{1em}{}
\titlespacing{\subsection}{0pt}{24pt}{6pt}

\titleformat{\subsubsection}{\raggedright\heiti\fontsize{13pt}{\baselineskip}\selectfont}{\thesubsubsection}{1em}{}
\titlespacing{\subsubsection}{0pt}{12pt}{6pt}

\lstset{
    language=Python,
    basicstyle=\ttfamily\fontsize{10.5pt}{16pt}\selectfont,
    keywordstyle=\color{blue},
    commentstyle=\color{green!60!black},
    stringstyle=\color{red},
    numbers=left,
    numberstyle=\tiny\color{gray},
    frame=single,
    breaklines=true,
    showstringspaces=false,
    tabsize=4,
    xleftmargin=2em,
    framexleftmargin=2em
}

\begin{document}

\begin{center}
{\heiti\fontsize{22pt}{\baselineskip}\selectfont\bfseries 强化学习第二次作业}
\end{center}

\vspace{12pt}

\begin{center}
{\fontsize{14pt}{\baselineskip}\selectfont 2312167\quad 刘振毫}
\end{center}

\vspace{24pt}

\begin{center}
{\heiti\fontsize{18pt}{\baselineskip}\selectfont\bfseries 摘要}
\end{center}

\vspace{18pt}

\setlength{\parindent}{2em}
\setlength{\parskip}{0pt}
\setstretch{1.0}

本文基于Gymnasium库开展强化学习实验研究，主要包含两个部分。第一部分选择CartPole-v1游戏环境，详细分析了其马尔可夫决策过程(MDP)模型的五个核心要素：状态空间、动作空间、状态转移概率、奖励函数和折扣因子，建立了完整的MDP数学描述。第二部分将经典的3臂赌博机问题扩展为4臂赌博机，新增摇臂期望奖励设为3.0。实验实现了四种探索策略：贪婪策略($\epsilon$-Greedy)、玻尔兹曼策略(Softmax)、UCB策略(Upper Confidence Bound)和汤普森采样(Thompson Sampling)。通过1000步训练、100次重复实验，系统比较了各策略在平均累积奖励和摇臂选择分布上的性能表现。实验结果表明，UCB策略和汤普森采样表现最优，最终平均累积奖励分别达到2.9475和2.9260，选择最优摇臂比例分别为96.1\%和94.1\%，能够有效平衡探索与利用。

\vspace{12pt}

\noindent\textbf{关键词：}强化学习；马尔可夫决策过程；多臂赌博机；探索与利用；UCB策略

\vspace{24pt}

\begin{center}
{\fontsize{18pt}{\baselineskip}\selectfont\bfseries Abstract}
\end{center}

\vspace{18pt}

This paper conducts reinforcement learning experiments based on the Gymnasium library, consisting of two main parts. The first part selects the CartPole-v1 game environment and provides a detailed analysis of its Markov Decision Process (MDP) model's five core elements: state space, action space, transition probability, reward function, and discount factor, establishing a complete mathematical description of the MDP. The second part extends the classic 3-armed bandit problem to a 4-armed bandit, with the newly added arm having an expected reward of 3.0. The experiment implements four exploration strategies: $\epsilon$-Greedy, Softmax (Boltzmann), UCB (Upper Confidence Bound), and Thompson Sampling. Through 1000-step training with 100 repeated experiments, the performance of each strategy in terms of average cumulative reward and arm selection distribution is systematically compared. Experimental results show that UCB and Thompson Sampling perform best, with final average cumulative rewards of 2.9475 and 2.9260 respectively, and optimal arm selection ratios of 96.1\% and 94.1\%, effectively balancing exploration and exploitation.

\vspace{12pt}

\noindent\textbf{Key Words:} Reinforcement Learning; Markov Decision Process; Multi-Armed Bandit; Exploration and Exploitation; UCB Strategy

\newpage

\section{实验目的}

\setlength{\parindent}{2em}
\setlength{\parskip}{0pt}
\setstretch{1.0}

本实验旨在实现以下目标：

\begin{enumerate}
    \item 学习使用Gymnasium库进行强化学习环境交互，深入理解马尔可夫决策过程(MDP)的基本要素，掌握状态空间、动作空间、转移概率、奖励函数和折扣因子的概念与建模方法。
    
    \item 将3臂赌博机扩展为4臂赌博机，新增摇臂期望奖励设为3.0，比较不同探索策略（贪婪策略、玻尔兹曼策略、UCB策略、汤普森采样）的性能表现。
    
    \item 分析训练次数与平均总回报的关系，深入理解强化学习中探索与利用的平衡问题。
\end{enumerate}

\section{实验环境}

本实验的软硬件环境配置如下：

\begin{itemize}
    \item 操作系统：Windows
    \item 编程语言：Python 3.12
    \item 主要依赖库：gymnasium 1.2.3, numpy 2.2.6, matplotlib

\end{itemize}

\section{任务一：Gym库游戏MDP建模}

\subsection{游戏选择}

本实验选择Gymnasium库中的CartPole-v1游戏。CartPole是强化学习领域的经典控制问题，目标是通过左右移动小车来保持杆子平衡。小车可以在一维轨道上自由移动，杆子一端固定在小车上，初始状态为直立。智能体需要通过施加左或右的力来控制小车移动，使杆子尽可能长时间保持直立不倒。

\subsection{MDP模型分析}

马尔可夫决策过程(MDP)由五元组$(S, A, P, R, \gamma)$组成，CartPole-v1的MDP模型详细分析如下：

\subsubsection{状态空间 S}

状态是一个4维连续向量，描述了小车和杆子的完整物理状态：

\begin{equation}
s = (x, \dot{x}, \theta, \dot{\theta})
\end{equation}

其中各分量的物理含义和取值范围如下：

\begin{itemize}
    \item $x \in [-4.8, 4.8]$：小车在轨道上的位置坐标
    \item $\dot{x} \in (-\infty, +\infty)$：小车的运动速度
    \item $\theta \in [-0.418, 0.418]$ rad（约$\pm 24°$）：杆子与垂直方向的夹角
    \item $\dot{\theta} \in (-\infty, +\infty)$：杆子的角速度
\end{itemize}

\subsubsection{动作空间 A}

CartPole-v1采用离散动作空间，包含2个动作：

\begin{itemize}
    \item $a=0$：向左推小车，施加负向力
    \item $a=1$：向右推小车，施加正向力
\end{itemize}

每个动作会对小车施加固定大小的力，力的方向由动作决定。

\subsubsection{状态转移概率 P}

状态转移概率$P(s'|s,a)$描述了在状态$s$下执行动作$a$后转移到状态$s'$的概率：

\begin{equation}
P(s'|s,a) = \delta(s' - f(s,a))
\end{equation}

其中$f(s,a)$是基于经典力学方程的确定性转移函数。CartPole的状态转移具有以下特点：

\begin{itemize}
    \item 转移是确定性的，不存在随机噪声
    \item 基于牛顿力学方程计算下一状态
    \item 终止条件：杆子角度超过$\pm 12°$（约0.209 rad）或小车位置超过$\pm 2.4$
\end{itemize}

\subsubsection{奖励函数 R}

奖励函数$R(s,a,s')$定义了智能体在状态$s$执行动作$a$后转移到$s'$时获得的即时奖励：

\begin{equation}
R(s,a,s') = 
\begin{cases}
+1 & \text{若游戏继续}\\
0 & \text{若游戏结束}
\end{cases}
\end{equation}

每保持一步平衡，智能体获得+1的奖励。游戏目标是最大化累积奖励，即尽可能长时间保持杆子平衡。CartPole-v1的最大步数为500步。

\subsubsection{折扣因子 $\gamma$}

折扣因子$\gamma \in [0,1]$决定了未来奖励的当前价值：

\begin{equation}
G_t = \sum_{k=0}^{\infty} \gamma^k R_{t+k+1}
\end{equation}

通常设为$\gamma = 0.99$或$\gamma = 1$。由于CartPole每步奖励相同（均为+1），折扣因子的选择对策略学习影响不大。

\subsection{代码实现与运行结果}

CartPole-v1环境的核心运行代码如下：

\begin{lstlisting}[caption={CartPole-v1环境运行代码}]
def task1_cartpole():
    try:
        import gymnasium as gym
    except ImportError:
        import gym
    
    env = gym.make('CartPole-v1', render_mode=None)
    
    observation, info = env.reset()
    total_reward = 0
    
    for step in range(100):
        action = env.action_space.sample()
        observation, reward, terminated, truncated, info = env.step(action)
        total_reward += reward
        
        if terminated or truncated:
            print(f"游戏结束，总步数: {step+1}, 总奖励: {total_reward}")
            break
    
    env.close()
    return total_reward
\end{lstlisting}

运行结果示例如下：

\begin{lstlisting}[caption={CartPole-v1运行示例输出}]
步骤1: 状态=['0.033', '0.161', '-0.003', '-0.300'], 动作=1, 奖励=1.0
步骤2: 状态=['0.037', '-0.034', '-0.009', '-0.008'], 动作=0, 奖励=1.0
步骤3: 状态=['0.036', '-0.229', '-0.009', '0.281'], 动作=0, 奖励=1.0
步骤4: 状态=['0.031', '-0.424', '-0.004', '0.571'], 动作=0, 奖励=1.0
步骤5: 状态=['0.023', '-0.619', '0.008', '0.863'], 动作=0, 奖励=1.0
游戏结束，总步数: 12, 总奖励: 12.0
\end{lstlisting}

\section{任务二：四臂赌博机实验}

\subsection{问题描述}

多臂赌博机(Multi-Armed Bandit, MAB)是强化学习的经典问题，体现了探索与利用的核心矛盾。本实验将3臂赌博机扩展为4臂赌博机，新增摇臂的期望奖励设为3.0。

\subsubsection{四臂赌博机参数设置}

四臂赌博机的参数设置如表1所示：

\begin{table}[H]
\centering
\caption{四臂赌博机参数设置}
\begin{tabular}{ccccc}
\toprule
摇臂编号 & 0 & 1 & 2 & 3（新增）\\
\midrule
期望奖励 $\mu$ & 1.0 & 2.0 & 1.5 & 3.0 \\
奖励分布 & $\mathcal{N}(1.0, 1)$ & $\mathcal{N}(2.0, 1)$ & $\mathcal{N}(1.5, 1)$ & $\mathcal{N}(3.0, 1)$ \\
\bottomrule
\end{tabular}
\end{table}

最优摇臂为摇臂3，期望奖励为3.0。每个摇臂的奖励服从正态分布，标准差均为1.0。

\subsection{策略实现}

本实验实现了四种探索策略，每种策略采用不同的方法平衡探索与利用。

\subsubsection{贪婪策略}

贪婪策略($\epsilon$-Greedy)是最简单的探索策略，以概率$\epsilon$随机探索，以概率$1-\epsilon$选择当前最优动作：

\begin{equation}
a_t = 
\begin{cases}
\arg\max_a Q_t(a) & \text{概率 } 1-\epsilon\\
\text{随机动作} & \text{概率 } \epsilon
\end{cases}
\end{equation}

其中$Q_t(a)$表示动作$a$在时刻$t$的估计价值，采用增量式更新：

\begin{equation}
Q_{t+1}(a) = Q_t(a) + \frac{1}{N_t(a)}[R_t - Q_t(a)]
\end{equation}

\subsubsection{玻尔兹曼策略}

玻尔兹曼策略(Softmax)根据动作值的Softmax分布选择动作，温度参数$T$控制探索程度：

\begin{equation}
P(a) = \frac{e^{Q(a)/T}}{\sum_{a'} e^{Q(a')/T}}
\end{equation}

温度$T$越高，选择各动作的概率越均匀（探索越多）；温度$T$越低，越倾向于选择当前最优动作（利用越多）。

\subsubsection{UCB策略}

UCB策略(Upper Confidence Bound)在动作值基础上增加不确定性bonus，实现探索与利用的平衡：

\begin{equation}
a_t = \arg\max_a \left[ Q_t(a) + c\sqrt{\frac{\ln t}{N_t(a)}} \right]
\end{equation}

其中$c$是探索参数，$N_t(a)$是动作$a$被选择的次数。该策略基于置信区间上界原则，被选择次数少的动作会获得更高的探索bonus。

\subsubsection{汤普森采样}

汤普森采样(Thompson Sampling)是一种贝叶斯方法，通过从后验分布采样来选择动作：

\begin{equation}
a_t = \arg\max_a \theta_t(a), \quad \theta_t(a) \sim \text{Beta}(\alpha_a, \beta_a)
\end{equation}

其中$\alpha_a$和$\beta_a$是Beta分布的参数，根据观测到的奖励进行更新。汤普森采样自然地实现了探索与利用的平衡。

\subsection{实验设置}

实验参数设置如下：

\begin{itemize}
    \item 每轮训练步数：1000步
    \item 重复次数：100次取平均
    \item 评估指标：平均累积奖励、各摇臂选中次数
    \item 贪婪策略参数：$\epsilon = 0.1$ 和 $\epsilon = 0.01$
    \item 玻尔兹曼策略参数：$T = 0.5$
    \item UCB策略参数：$c = 2.0$
\end{itemize}

\subsection{实验结果}

\subsubsection{各策略最终平均累积奖励}

各策略最终平均累积奖励排名如表2所示：

\begin{table}[H]
\centering
\caption{各策略最终平均累积奖励排名}
\begin{tabular}{ccc}
\toprule
排名 & 策略 & 最终平均累积奖励 \\
\midrule
1 & UCB策略(c=2) & 2.9475 \\
2 & 汤普森采样 & 2.9260 \\
3 & 贪婪策略($\epsilon$=0.1) & 2.8431 \\
4 & 玻尔兹曼策略(T=0.5) & 2.8021 \\
5 & 贪婪策略($\epsilon$=0.01) & 2.5886 \\
\bottomrule
\end{tabular}
\end{table}

\subsubsection{各策略摇臂选中分布}

各策略下每个摇臂的选中次数分布如表3所示：

\begin{table}[H]
\centering
\caption{各策略下每个摇臂的选中次数分布}
\begin{tabular}{ccccc}
\toprule
策略 & 摇臂0 ($\mu$=1.0) & 摇臂1 ($\mu$=2.0) & 摇臂2 ($\mu$=1.5) & 摇臂3 ($\mu$=3.0) \\
\midrule
贪婪策略($\epsilon$=0.1) & 4.7\% & 3.8\% & 2.8\% & 88.7\% \\
贪婪策略($\epsilon$=0.01) & 3.5\% & 19.0\% & 10.0\% & 67.5\% \\
玻尔兹曼策略(T=0.5) & 1.6\% & 10.9\% & 4.1\% & 83.4\% \\
UCB策略(c=2) & 0.7\% & 2.1\% & 1.1\% & 96.1\% \\
汤普森采样 & 1.0\% & 3.3\% & 1.6\% & 94.1\% \\
\bottomrule
\end{tabular}
\end{table}

\subsubsection{实验结果图}

实验结果如图1所示，包含四个子图：训练次数与平均累积奖励关系、各摇臂选中次数、收敛过程对比和最终性能对比。

\begin{figure}[H]
\centering
\includegraphics[width=0.95\textwidth]{bandit_results.png}
\caption{四臂赌博机实验结果}
\end{figure}

\section{结果分析与讨论}

\subsection{策略性能比较}

从实验结果可以得出以下结论：

\textbf{UCB策略表现最优}：最终平均累积奖励达到2.9475，选择最优摇臂比例高达96.1\%。UCB通过不确定性bonus有效平衡了探索与利用，能够快速识别最优摇臂。其理论保证使其在多臂赌博机问题上具有优异性能。

\textbf{汤普森采样紧随其后}：最终奖励2.9260，选择最优摇臂比例94.1\%。作为贝叶斯方法，汤普森采样通过后验分布采样自然地实现了探索与利用的平衡，无需手动调节探索参数。

\textbf{贪婪策略受$\epsilon$影响显著}：$\epsilon$=0.1时表现较好（2.8431），但仍有一定探索损失；$\epsilon$=0.01时表现最差（2.5886），探索不足导致容易陷入次优解。这表明固定探索率难以适应不同阶段的需求。

\textbf{玻尔兹曼策略表现中等}：温度参数T=0.5时，表现介于两种贪婪策略之间。温度参数的选择对性能有较大影响，需要根据具体问题调参。

\subsection{探索与利用的平衡}

从实验结果可以看出探索与利用平衡的重要性：

\begin{itemize}
    \item UCB和汤普森采样能够自适应地调整探索程度，在初期充分探索后快速收敛到最优策略
    \item 固定探索率的贪婪策略无法自适应调整，$\epsilon$过大造成不必要的探索损失，$\epsilon$过小则探索不足
    \item 玻尔兹曼策略的温度参数同样需要手动调节，难以达到最优平衡
\end{itemize}

\subsection{新增摇臂的影响}

新增的摇臂3期望奖励为3.0，成为最优摇臂。实验表明：

\begin{itemize}
    \item UCB和汤普森采样能够快速识别并利用新增的最优摇臂
    \item 贪婪策略($\epsilon$=0.01)由于探索不足，有较高概率错过最优摇臂，67.5\%的选中率远低于其他策略
\end{itemize}

\section{结论}

本实验得出以下结论：

\begin{enumerate}
    \item 成功使用Gymnasium库运行CartPole-v1游戏，并建立了完整的MDP模型$(S, A, P, R, \gamma)$，深入理解了强化学习环境的基本要素。
    
    \item 将3臂赌博机成功扩展为4臂赌博机，新增摇臂期望奖励设为3.0，成为最优摇臂。
    
    \item 实现了四种探索策略：贪婪策略、玻尔兹曼策略、UCB策略和汤普森采样，并通过实验系统比较了各策略的性能。
    
    \item 实验结果表明，UCB策略和汤普森采样在四臂赌博机问题上表现最优，能够有效平衡探索与利用，最终平均累积奖励分别达到2.9475和2.9260。
    
    \item 固定探索率的贪婪策略性能受参数选择影响较大，需要根据具体问题调参，$\epsilon$=0.01时表现最差。
\end{enumerate}

\newpage

\section*{附录}

\addcontentsline{toc}{section}{附录}

\subsection*{完整代码}

\begin{lstlisting}[caption={强化学习实验完整代码}]
import numpy as np
import matplotlib.pyplot as plt
from matplotlib import rcParams

rcParams['font.sans-serif'] = ['SimHei']
rcParams['axes.unicode_minus'] = False

#Gym库游戏 - CartPole-v1
def task1_cartpole():

    try:
        import gymnasium as gym
    except ImportError:
        import gym
    
    env = gym.make('CartPole-v1', render_mode=None)
   
    print("\n【运行演示】")
    observation, info = env.reset()
    total_reward = 0
    
    for step in range(100):
        action = env.action_space.sample()
        observation, reward, terminated, truncated, info = env.step(action)
        total_reward += reward
        
        if step < 5:
            print(f"  步骤{step+1}: 状态={[f'{x:.3f}' for x in observation]}, 动作={action}, 奖励={reward}")
        
        if terminated or truncated:
            print(f"  游戏结束，总步数: {step+1}, 总奖励: {total_reward}")
            break
    
    env.close()
    return total_reward


# 任务2：4臂赌博机

class FourArmedBandit:
    def __init__(self):
        self.n_arms = 4
        self.true_means = [1.0, 2.0, 1.5, 3.0]
        self.std = 1.0
        self.best_arm = np.argmax(self.true_means)
    
    def pull(self, arm):
        return np.random.normal(self.true_means[arm], self.std)
    
    def get_optimal_reward(self):
        return self.true_means[self.best_arm]


class EpsilonGreedyAgent:
    def __init__(self, n_arms, epsilon=0.1):
        self.n_arms = n_arms
        self.epsilon = epsilon
        self.counts = np.zeros(n_arms)
        self.values = np.zeros(n_arms)
    
    def select_action(self):
        if np.random.random() < self.epsilon:
            return np.random.randint(self.n_arms)
        else:
            max_value = np.max(self.values)
            candidates = np.where(self.values == max_value)[0]
            return np.random.choice(candidates)
    
    def update(self, arm, reward):
        self.counts[arm] += 1
        self.values[arm] += (reward - self.values[arm]) / self.counts[arm]


class BoltzmannAgent:
    def __init__(self, n_arms, temperature=1.0):
        self.n_arms = n_arms
        self.temperature = temperature
        self.counts = np.zeros(n_arms)
        self.values = np.zeros(n_arms)
    
    def select_action(self):
        if self.temperature == 0:
            return np.argmax(self.values)
        
        exp_values = np.exp(self.values / self.temperature)
        probs = exp_values / np.sum(exp_values)
        return np.random.choice(self.n_arms, p=probs)
    
    def update(self, arm, reward):
        self.counts[arm] += 1
        self.values[arm] += (reward - self.values[arm]) / self.counts[arm]


class UCBAgent:
    def __init__(self, n_arms, c=2.0):
        self.n_arms = n_arms
        self.c = c
        self.counts = np.zeros(n_arms)
        self.values = np.zeros(n_arms)
        self.total_count = 0
    
    def select_action(self):
        if self.total_count < self.n_arms:
            return self.total_count
        
        ucb_values = np.zeros(self.n_arms)
        for arm in range(self.n_arms):
            if self.counts[arm] == 0:
                ucb_values[arm] = float('inf')
            else:
                bonus = self.c * np.sqrt(np.log(self.total_count) / self.counts[arm])
                ucb_values[arm] = self.values[arm] + bonus
        
        return np.argmax(ucb_values)
    
    def update(self, arm, reward):
        self.counts[arm] += 1
        self.total_count += 1
        self.values[arm] += (reward - self.values[arm]) / self.counts[arm]


class ThompsonSamplingAgent:
    def __init__(self, n_arms):
        self.n_arms = n_arms
        self.alpha = np.ones(n_arms)
        self.beta = np.ones(n_arms)
        self.counts = np.zeros(n_arms)
        self.values = np.zeros(n_arms)
    
    def select_action(self):
        samples = np.random.beta(self.alpha, self.beta)
        return np.argmax(samples)
    
    def update(self, arm, reward):
        self.counts[arm] += 1
        self.values[arm] += (reward - self.values[arm]) / self.counts[arm]
        
        normalized_reward = (reward + 3) / 6
        normalized_reward = np.clip(normalized_reward, 0, 1)
        
        self.alpha[arm] += normalized_reward
        self.beta[arm] += (1 - normalized_reward)


def run_experiment(agent_class, agent_params, bandit, n_steps=1000, n_runs=100):
    all_rewards = np.zeros((n_runs, n_steps))
    all_arm_counts = np.zeros((n_runs, bandit.n_arms))
    
    for run in range(n_runs):
        agent = agent_class(**agent_params)
        for step in range(n_steps):
            arm = agent.select_action()
            reward = bandit.pull(arm)
            agent.update(arm, reward)
            all_rewards[run, step] = reward
        
        all_arm_counts[run] = agent.counts
    
    avg_rewards = np.mean(all_rewards, axis=0)
    cumulative_avg_rewards = np.cumsum(avg_rewards) / np.arange(1, n_steps + 1)
    avg_arm_counts = np.mean(all_arm_counts, axis=0)
    
    return cumulative_avg_rewards, avg_arm_counts


def task2_bandit():
    print("任务2：4臂赌博机实验")
    print("【4臂赌博机设置】")
    bandit = FourArmedBandit()
    print(f"摇臂数量: {bandit.n_arms}")
    print(f"各摇臂真实期望奖励: {bandit.true_means}")
    print(f"奖励标准差: {bandit.std}")
    print(f"最优摇臂: 摇臂{bandit.best_arm} (期望={bandit.true_means[bandit.best_arm]})")
    
    n_steps = 1000
    n_runs = 100
    
    print(f"\n实验设置: 每轮{n_steps}步, 重复{n_runs}次取平均")
    
    agents_config = {
        '贪婪策略(epsilon=0.1)': (EpsilonGreedyAgent, {'n_arms': 4, 'epsilon': 0.1}),
        '贪婪策略(epsilon=0.01)': (EpsilonGreedyAgent, {'n_arms': 4, 'epsilon': 0.01}),
        '玻尔兹曼策略(T=0.5)': (BoltzmannAgent, {'n_arms': 4, 'temperature': 0.5}),
        'UCB策略(c=2)': (UCBAgent, {'n_arms': 4, 'c': 2.0}),
        '汤普森采样': (ThompsonSamplingAgent, {'n_arms': 4}),
    }
    
    results = {}
    for name, (agent_class, params) in agents_config.items():
        print(f"  运行 {name}...")
        cumulative_avg, arm_counts = run_experiment(agent_class, params, bandit, n_steps, n_runs)
        results[name] = {
            'cumulative_avg_rewards': cumulative_avg,
            'arm_counts': arm_counts
        }
    
    fig, axes = plt.subplots(2, 2, figsize=(14, 10))
    
    ax1 = axes[0, 0]
    for name, data in results.items():
        ax1.plot(data['cumulative_avg_rewards'], label=name, linewidth=1.5)
    ax1.axhline(y=bandit.get_optimal_reward(), color='red', linestyle='--', 
                label=f'最优期望奖励 ({bandit.get_optimal_reward()})', linewidth=2)
    ax1.set_xlabel('训练步数', fontsize=12)
    ax1.set_ylabel('平均累积奖励', fontsize=12)
    ax1.set_title('各策略训练次数与平均累积奖励的关系', fontsize=14)
    ax1.legend(loc='lower right', fontsize=9)
    ax1.grid(True, alpha=0.3)
    
    ax2 = axes[0, 1]
    arm_names = [f'摇臂{i}\n(mu={bandit.true_means[i]})' for i in range(4)]
    x = np.arange(4)
    width = 0.15
    
    for i, (name, data) in enumerate(results.items()):
        bars = ax2.bar(x + i * width, data['arm_counts'], width, label=name)
    
    ax2.set_xlabel('摇臂', fontsize=12)
    ax2.set_ylabel('平均选中次数', fontsize=12)
    ax2.set_title('各策略下每个摇臂的选中次数', fontsize=14)
    ax2.set_xticks(x + width * 2)
    ax2.set_xticklabels(arm_names)
    ax2.legend(loc='upper left', fontsize=8)
    ax2.grid(True, alpha=0.3, axis='y')
    
    ax3 = axes[1, 0]
    for name, data in results.items():
        cumulative = data['cumulative_avg_rewards']
        ax3.plot(cumulative, label=name, linewidth=1.5, alpha=0.7)
    ax3.axhline(y=bandit.get_optimal_reward(), color='red', linestyle='--', 
                label='最优期望奖励', linewidth=2)
    ax3.set_xlabel('训练步数', fontsize=12)
    ax3.set_ylabel('平均累积奖励', fontsize=12)
    ax3.set_title('收敛过程对比', fontsize=14)
    ax3.legend(loc='lower right', fontsize=9)
    ax3.grid(True, alpha=0.3)
    ax3.set_xlim(0, 200)
    
    ax4 = axes[1, 1]
    strategies = list(results.keys())
    final_rewards = [results[s]['cumulative_avg_rewards'][-1] for s in strategies]
    colors = plt.cm.Set3(np.linspace(0, 1, len(strategies)))
    
    bars = ax4.barh(strategies, final_rewards, color=colors)
    ax4.axvline(x=bandit.get_optimal_reward(), color='red', linestyle='--', 
                label=f'最优期望奖励', linewidth=2)
    ax4.set_xlabel('最终平均累积奖励', fontsize=12)
    ax4.set_title('各策略最终性能对比', fontsize=14)
    ax4.legend(loc='lower right', fontsize=9)
    ax4.grid(True, alpha=0.3, axis='x')
    
    for bar, reward in zip(bars, final_rewards):
        ax4.text(bar.get_width() + 0.02, bar.get_y() + bar.get_height()/2, 
                f'{reward:.3f}', va='center', fontsize=9)
    
    plt.tight_layout()
    plt.savefig('bandit_results.png', dpi=150, bbox_inches='tight')
    plt.show()
    
    print("【实验结果总结】")
    print(f"\n各摇臂真实期望: {bandit.true_means}")
    print(f"最优摇臂: 摇臂{bandit.best_arm} (期望={bandit.true_means[bandit.best_arm]})")
    print("\n各策略最终平均累积奖励:")
    for name, data in sorted(results.items(), key=lambda x: x[1]['cumulative_avg_rewards'][-1], reverse=True):
        final_reward = data['cumulative_avg_rewards'][-1]
        print(f"  {name}: {final_reward:.4f}")
    
    print("\n各策略摇臂选中分布:")
    for name, data in results.items():
        counts = data['arm_counts']
        total = sum(counts)
        percentages = [c/total*100 for c in counts]
        print(f"  {name}:")
        for i, (cnt, pct) in enumerate(zip(counts, percentages)):
            print(f"    摇臂{i}: {cnt:.1f}次 ({pct:.1f}%)")
    
    print("\n图表已保存至: bandit_results.png")



if __name__ == "__main__":   
    task1_cartpole()
    task2_bandit()
    
    print("\n所有任务完成！")
\end{lstlisting}

\end{document}
